% ============================================================
%  lang/pl.tex --- każdy widoczny dla czytelnika napis w wydaniu polskim
%
%  Nic w chapters/pl/ nie może wpisywać na sztywno słowa, które składa też
%  makro. Jeśli potrzebna jest etykieta, której tu nie ma, dodaj ją tutaj I w
%  lang/en.tex; CI przerywa build, gdy oba pliki nie definiują tego samego
%  zbioru makr (C3).
% ============================================================

% ---------- Tytuły części ----------
\newcommand{\lblPartI}{Mechanizm zegara: jak matematyka opisuje świat}
\newcommand{\lblPartII}{Gdzie kończy się przewidywanie}
\newcommand{\lblPartIII}{Kształty chaosu}
\newcommand{\lblPartIV}{Chaos w świecie}
\newcommand{\lblPartV}{Życie z chaosem}
\newcommand{\lblPartAppendices}{Dodatki}

% ---------- Tytuły ramek ----------
\newcommand{\lblNote}{Uwaga}
\newcommand{\lblWarning}{Ostrożnie}
\newcommand{\lblTrap}{Pułapka}
\newcommand{\lblWorld}{Gdzie to spotkasz}
\newcommand{\lblRigour}{Czego tu nie dowodzimy}
\newcommand{\lblProject}{chaoslab}
\newcommand{\lblVerify}{Uruchom, zanim zaufasz}
\newcommand{\lblLab}{Laboratorium}

% ---------- Metoda ----------
\newcommand{\lblThink}{Zatrzymaj się i pomyśl}
\newcommand{\lblThinkCue}{Odpowiedz, zanim pójdziesz dalej.}
\newcommand{\lblAnswerMark}{Odpowiedź.}
\newcommand{\lblOutcomes}{Czego się nauczysz}
\newcommand{\lblOutcomesLead}{Po przerobieniu tego rozdziału będziesz umieć:}
\newcommand{\lblSummary}{Podsumowanie}
\newcommand{\lblCanYou}{Czy potrafisz?}
\newcommand{\lblCanYouIntro}{%
  Oceń się przy każdym punkcie od 1 (nie) do 5 (tak, samodzielnie). Wszystko
  poniżej 4 nie jest powodem do rozmyślań: wróć do podrozdziału, który tego
  uczył, i jeszcze raz odpowiedz na jego pytania, zanim pójdziesz dalej.}
\newcommand{\lblNoOutcomes}{W tym rozdziale nie zadeklarowano celów.}
\newcommand{\lblCheckpoint}{Punkt kontrolny}
\newcommand{\lblCheckpointIntro}{%
  Odpowiedz na te pytania na papierze, zanim zajrzysz do dodatku~A.
  Sprawdzają to, czego rozdział uczył, w tej kolejności, w jakiej uczył.}
\newcommand{\lblAnswersIntro}{%
  Odpowiedzi do punktów kontrolnych wszystkich rozdziałów, w kolejności
  rozdziałów. Każdą zapisano obok jej pytania, a build zebrał je tutaj, więc
  odpowiedź nie może odpłynąć od pytania, na które odpowiada.}
\newcommand{\lblNoAnswers}{Nie zapisano żadnych odpowiedzi. To jest build częściowy.}

% ---------- Stopka laboratorium ----------
\newcommand{\lblLabStarter}{Plik startowy}
\newcommand{\lblLabTest}{Test}
\newcommand{\lblLabRun}{Uruchom}

% ---------- Znaczniki stanu builda ----------
\newcommand{\lblNotWritten}{JESZCZE NIENAPISANE}
\newcommand{\lblNotRendered}{DIAGRAM NIEWYRENDEROWANY --- uruchom \texttt{make diagrams}}
\newcommand{\lblPlotNotRendered}{WYKRES NIEWYRENDEROWANY --- uruchom \texttt{make plots}}
\newcommand{\lblTranscriptMissing}{TRANSKRYPT NIEWYGENEROWANY --- uruchom \texttt{make numbers}}
\newcommand{\lblValueMissing}{brak wartości}
\newcommand{\lblDiagramManifest}{Manifest diagramów}
\newcommand{\lblPlotManifest}{Manifest wykresów}
\newcommand{\lblLabManifest}{Manifest laboratoriów}

% ---------- Żywe paginy frontmatter ----------
\newcommand{\lblHowToUse}{Jak korzystać z tej książki}
\newcommand{\lblIntroduction}{Wprowadzenie}

% Myślnik interpunkcyjny, pisany \dash{} ze spacją po obu stronach: polskie
% wydanie składa półpauzę, angielskie pauzę.
\newcommand{\dash}{--}

% ---------- Odsyłacze w indeksie ----------
\newcommand{\lblIndexSee}{zob.}
\newcommand{\lblIndexSeeAlso}{zob.\ też}

% Data ostatniego sprawdzenia faktów i przypiętych wersji.
\newcommand{\pinnedon}{wrzesień 2026}

% ---------- Metadane PDF ----------
\newcommand{\lblPdfSubject}{Jak matematyka opisuje świat i gdzie przestaje go przewidywać: teoria chaosu od szkolnej arytmetyki}
\newcommand{\lblPdfKeywords}{teoria chaosu, układy dynamiczne, odwzorowanie logistyczne, atraktor Lorenza, fraktale, wykładnik Lapunowa, przewidywanie, Python}
\newcommand{\lblPdfLang}{pl-PL}
